\SixthArticle{16}
Wir werden diese Untersuchung allgemeiner auf eine unendliche Reihe $M,\allowbreak M',\allowbreak M'',\allowbreak M'''$ usw. anwenden, die so gebildet ist, dass die Quotienten $\frac{M'}{M},\frac{M''}{M'},\frac{M'''}{M''}$ usw. jeweils die Werte des Bruchs
\[
\frac{t^{\lambda}+At^{\lambda-1}+Bt^{\lambda-2}+Ct^{\lambda-3}+\mathrm{usw.}}
{t^{\lambda}+at^{\lambda-1}+bt^{\lambda-2}+ct^{\lambda-3}+\mathrm{usw.}}
\]
für $t=m$, $t=m+1$, $t=m+2$ usw. sind. Der Kürze halber werden wir den Zähler dieses Bruchs mit $P$, den Nenner mit $p$ bezeichnen. Außerdem werden wir voraussetzen, dass $P,p$ nicht identisch sind, also die Differenzen $A-a,B-b,C-c$ usw. nicht sämtlich zugleich verschwinden.

I. Ist von den Differenzen $A-a,B-b,C-c$ usw. die erste, die nicht verschwindet, positiv, so lässt sich eine Grenze $l$ angeben, nach deren Überschreitung durch den Wert von $t$ die Werte der Funktionen $P$ und $p$ gewiss stets positiv werden und $P>p$ gilt. Dies tritt offenbar ein, wenn als $l$ die größte reelle Wurzel der Gleichung $p(P-p)=0$ genommen wird; hat diese Gleichung aber überhaupt keine reellen Wurzeln, so gilt jene Eigenschaft für alle Werte von $t$. Daher sind in der Reihe $\frac{M'}{M},\frac{M''}{M'},\frac{M'''}{M''}$ usw. wenigstens nach einem gewissen Abschnitt, wenn nicht von Anfang an, alle Glieder positiv und größer als die Einheit. Wird also keines $=0$ oder unendlich groß, so ist klar, dass
\SourceLeaf{149}
\textit{die Reihe $M,M',M'',M'''$ usw., wenn nicht von Anfang an, so doch nach einem gewissen Abschnitt, lauter Glieder desselben Vorzeichens besitzt, die stetig wachsen.}

Ist in gleicher Weise von den Differenzen $A-a,B-b,C-c$ usw. die erste, die nicht verschwindet, negativ, so besitzt die Reihe $M,M',M'',M'''$ usw., wenn nicht von Anfang an, so doch nach einem gewissen Abschnitt, lauter Glieder desselben Vorzeichens, die stetig abnehmen.

II. Sind nun die Koeffizienten $A,a$ ungleich, so werden die Glieder der Reihe $M,M',M'',M'''$ usw. je nachdem, ob die Differenz $A-a$ positiv oder negativ ist, über alle Grenzen wachsen oder gegen Null abnehmen. Wir beweisen dies folgendermaßen. Ist $A-a$ eine positive Größe, so nehme man eine ganze Zahl $h$ so, dass $h(A-a)>1$ wird, und setze $\frac{M^h}{m}=N$, $\frac{M^{\prime h}}{m+1}=N'$, $\frac{M^{\prime\prime h}}{m+2}=N''$, $\frac{M^{\prime\prime\prime h}}{m+3}=N'''$ usw. sowie $tP^h=Q$, $(t+1)p^h=q$. Dann sind offenbar $\frac{N'}{N},\frac{N''}{N'},\frac{N'''}{N''}$ usw. die Werte des Bruchs $\frac{Q}{q}$, wenn $t=m$, $t=m+1$, $t=m+2$ usw. gesetzt wird; $Q,q$ selbst sind dagegen algebraische Funktionen folgender Form:
\[
\begin{array}{l}
Q=t^{\lambda h+1}+hAt^{\lambda h}+\mathrm{usw.}\\
q=t^{\lambda h+1}+(ha+1)t^{\lambda h}+\mathrm{usw.}
\end{array}
\]
Da nach Voraussetzung die Differenz $hA-(ha+1)$ eine positive Größe ist, werden die Glieder der Reihe $N,N',N'',N'''$ usw., wenn nicht von Anfang an, so doch nach einem gewissen Abschnitt, stetig wachsen (nach \textsc{i}). Daher müssen die Glieder der Reihe $mN,(m+1)N',(m+2)N''$, $(m+3)N'''$ usw. über alle Grenzen wachsen, und folglich auch die Glieder der Reihe $M,M',M'',M'''$ usw., deren Potenzen mit dem Exponenten $h$ jenen Gliedern gleich sind. Q.~E.~P.

Ist $A-a$ eine negative Größe, so muss man eine ganze Zahl $h$ so wählen, dass $h(a-A)$ größer als $1$ wird. Dann werden durch ähnliche Schlüsse die Glieder der Reihe
\[
mM^h,\quad (m+1)M^{\prime h},\quad (m+2)M^{\prime\prime h},\quad (m+3)M^{\prime\prime\prime h}\quad\mathrm{usw.}
\]
nach einem gewissen Abschnitt stetig abnehmen. Daher müssen die Glieder der Reihe $M^h$, $M^{\prime h},M^{\prime\prime h}$ usw. und somit auch die Glieder von $M,\allowbreak M',\allowbreak M'',\allowbreak M'''$ usw. gegen Null abnehmen. Q.~E.~S.

III. Sind dagegen die ersten Koeffizienten $A,a$ gleich, so werden sich die Glieder der Reihe $M,\allowbreak M',\allowbreak M'',\allowbreak M'''$ usw. stetig einer endlichen Grenze nähern. Wir beweisen dies folgendermaßen. Wir nehmen zunächst an, dass die Glieder der Reihe nach einem gewissen Abschnitt stetig wachsen, dass also von den Differenzen $B-b,C-c$ usw. die erste, die nicht verschwindet,
\SourceLeaf{150}
positiv ist. Es sei $h$ eine ganze Zahl, für die $h+b-B$ eine positive Größe wird, und man setze
\[
M\left(\frac{m}{m-1}\right)^h=N,\quad
M'\left(\frac{m+1}{m}\right)^h=N',\quad
M''\left(\frac{m+2}{m+1}\right)^h=N''\quad\mathrm{usw.}
\]
und $(tt-1)^hP=Q,\ t^{2h}p=q$, sodass $\frac{N'}{N},\frac{N''}{N'}$ usw. die Werte des Bruchs $\frac{Q}{q}$ sind, wenn $t=m,\ t=m+1$ usw. gesetzt wird. Da nun
\[
\begin{array}{l}
Q=t^{\lambda+2h}+At^{\lambda+2h-1}+(B-h)t^{\lambda+2h-2}\quad\mathrm{usw.}\\
q=t^{\lambda+2h}+At^{\lambda+2h-1}+bt^{\lambda+2h-2}\quad\mathrm{usw.}
\end{array}
\]
gilt und nach Voraussetzung $B-h-b$ eine negative Größe ist, werden die Glieder der Reihe $N,N',N'',N'''$ usw. wenigstens nach einem gewissen Abschnitt stetig abnehmen. Daher können die Glieder der Reihe $M,M',M'',M'''$ usw., die jeweils stets kleiner als jene sind, während ihres stetigen Wachstums sich nur einer endlichen Grenze nähern. Q.~E.~P.

Nehmen die Glieder der Reihe $M,M',M'',M'''$ usw. nach einem gewissen Abschnitt stetig ab, so muss für $h$ eine ganze Zahl gewählt werden, für die $h+B-b$ eine positive Größe ist. Durch ganz ähnliche Schlüsse wird sich zeigen lassen, dass die Glieder der Reihe
\[
M\left(\frac{m-1}{m}\right)^h,\quad
M'\left(\frac{m}{m+1}\right)^h,\quad
M''\left(\frac{m+1}{m+2}\right)^h\quad\mathrm{usw.}
\]
nach einem gewissen Abschnitt stetig wachsen. Daher können die Glieder der Reihe $M,M',M''$ usw., die jeweils stets größer als jene sind, während ihrer stetigen Abnahme notwendig nur gegen eine endliche Grenze abnehmen. Q.~E.~S.

IV. Was schließlich die \textit{Summe} der Reihe mit den Gliedern $M,M',$ $M'',M'''$ usw. in dem Fall betrifft, dass diese gegen Null abnehmen, so nehmen wir zunächst an, dass $A-a$ zwischen $0$ und $-1$ liegt, dass also $A+1-a$ entweder eine positive Größe oder $=0$ ist. Es sei $h$ eine positive ganze Zahl, und zwar beliebig in dem Fall, dass $A+1-a$ eine positive Größe ist, oder so gewählt, dass sie die Größe $h+m+A+B-b$ positiv macht, wenn $A+1-a=0$ ist. Dann ist
\[
\begin{array}{l}
P(t-(m+h-1))=t^{\lambda+1}+(A+1-m-h)t^\lambda\\\qquad{}+(B-A(m+h-1))t^{\lambda-1}\quad\mathrm{usw.}\\
p(t-(m+h))=t^{\lambda+1}+(a-m-h)t^\lambda\\\qquad{}+(b-a(m+h))t^{\lambda-1}\quad\mathrm{usw.}
\end{array}
\]
wobei entweder $A+1-m-h-(a-m-h)$ eine positive Größe sein wird oder, wenn diese $=0$ wird, wenigstens $B-A(m+h-1)-(b-a(m+h))$ positiv sein wird. Daher lässt sich (nach \textsc{i}) für die Größe $t$ ein Wert $l$ angeben, nach dessen Überschreitung die Werte des Bruchs $\frac{P(t-(m+h-1))}{p(t-(m+h))}$ stets positiv und größer als die Einheit werden. Es sei $n$ eine ganze Zahl,
\SourceLeaf{151}
die größer als $l$ und zugleich größer als $h$ ist, und die Glieder der Reihe $M,M',M'',M'''$ usw., die den Werten $t=m+n,\ t=m+n+1,\ t=m+n+2$ usw. entsprechen, seien $N,N',N'',N'''$ usw. Dann sind
\[
\frac{(n+1-h)N'}{(n-h)N},\quad
\frac{(n+2-h)N''}{(n+1-h)N'},\quad
\frac{(n+3-h)N'''}{(n+2-h)N''}\quad\mathrm{usw.}
\]
positive Größen, die größer als die Einheit sind; daher ist
\[
N'>\frac{(n-h)N}{n+1-h},\quad
N''>\frac{(n-h)N}{n+2-h},\quad
N'''>\frac{(n-h)N}{n+3-h}\quad\mathrm{usw.}
\]
und somit ist die Summe der Reihe $N+N'+N''+N'''+$ usw. größer als die Summe der Reihe
\[
(n-h)N\left(\frac{1}{n-h}+\frac{1}{n+1-h}+\frac{1}{n+2-h}+\frac{1}{n+3-h}+\mathrm{usw.}\right)
\]
wie viele Glieder auch zusammengefasst werden. Die letztere Reihe überschreitet aber alle Grenzen, wenn die Zahl der Glieder ins Unendliche wächst, da die Summe der Reihe $1+\frac12+\frac13+\frac14+$ usw., die bekanntlich unendlich ist, auch dann unendlich bleibt, wenn die anfänglichen Glieder $1+\frac12+\frac13+\mathrm{usw.}$ $+\frac{1}{n-1-h}$ weggelassen werden. Daher wächst die Summe der Reihe $N+N'+N''+N'''+$ usw. und somit auch die Summe von $M+M'+M''+M'''+$ usw., deren Teil sie ist, über alle Grenzen.

V. Ist dagegen $A-a$ eine negative Größe, deren absoluter Wert größer als die Einheit ist, so wird die Summe der ins Unendliche fortgesetzten Reihe $M+M'+M''+M'''+$ usw. gewiss endlich sein. Es sei nämlich $h$ eine positive Größe, die kleiner als $a-A-1$ ist. Dann wird sich durch ähnliche Schlüsse zeigen lassen, dass sich ein Wert $l$ der Größe $t$ angeben lässt, über den hinaus der Bruch $\frac{Pt}{p(t-h-1)}$ stets positive Werte kleiner als die Einheit annimmt. Wird nun für $n$ eine ganze Zahl genommen, die größer als $l,m,h+1$ ist, und werden die Glieder der Reihe $M,M',M'',M'''$ usw., die den Werten $t=n,\ t=n+1,\ t=n+2$ usw. entsprechen, mit $N,N',N''$ usw. bezeichnet, so ist
\[
N'<\frac{n-h-1}{n}\cdot N,\quad
N''<\frac{(n-h-1)(n-h)}{n(n+1)}\cdot N'\quad\mathrm{usw.}
\]
und somit ist die Summe der Reihe $N+N'+N''+$ usw., wie viele Glieder auch zusammengefasst werden, kleiner als das Produkt aus $N$ und der Summe ebenso vieler Glieder der Reihe
\[
1+\frac{n-h-1}{n}+\frac{(n-h-1)(n-h)}{n(n+1)}+
\frac{(n-h-1)(n-h)(n-h+1)}{n(n+1)(n+2)}\quad\mathrm{usw.}
\]
Deren Summe lässt sich für jede Anzahl von Gliedern leicht angeben; es ist nämlich
\[
\begin{array}{lr}
\mathrm{erstes\ Glied}&=\displaystyle\frac{n-1}{h}-\frac{n-h-1}{h}
\end{array}
\]
\SourceLeaf{152}
\[
\begin{array}{@{}l@{\;}l@{}}
\text{Summe zweier Glieder}&=\displaystyle\frac{n-1}{h}-\frac{(n-h-1)(n-h)}{hn}\\[7pt]
\text{Summe dreier Glieder}&=\displaystyle\frac{n-1}{h}-\frac{(n-h-1)(n-h)(n-h+1)}{hn(n+1)}\quad\text{usw.}
\end{array}
\]
Da der zweite Teil (nach II) eine gegen Null abnehmende Reihe bildet, muss jene ins Unendliche fortgesetzte Summe $=\frac{n-1}{h}$ gesetzt werden. Daher bleibt $N+N'+N''$ usw. ins Unendliche fortgesetzt stets kleiner als $\frac{N(n-1)}{h}$, und folglich wird $M+M'+M''$ usw. gewiss gegen eine endliche Summe konvergieren. Q.~E.~D.

\Needspace{8\baselineskip}
VI. Um das, was wir allgemein über die Reihe $M,M',M''$ usw. bewiesen haben, auf die Koeffizienten der Potenzen $x^m,x^{m+1},x^{m+2}$ usw. in der Reihe $F(\alpha,\beta,\gamma,x)$ anzuwenden, wird man $\lambda=2$, $A=\alpha+\beta$, $B=\alpha\beta$, $a=\gamma+1$, $b=\gamma$ setzen müssen. Daraus ergeben sich unmittelbar die fünf im vorigen Artikel aufgestellten Behauptungen.

